\scanpage{25}
\mspage{16}
% archive sheet number: 25
% lines 1--3
La condition que $M$ est \inserted{pré}polyédral \deleted{\unsure{est ici}} équivaut à dire que
\[
X\leq Y\Longrightarrow X\ll_{\mathrm{pol}}Y
\]
% lines 4--6
ce qui est nécessaire pour que
\[
(M,\leq,\ll_{\mathrm{pol}},\mid\circ\mid)
\]
soit encore un magasin. Est-ce suffisant ?

% lines 7--9
1°) Mag 1\quad Il reste à vérifier que $\ll_{\mathrm{pol}}$ est une relation d'ordre, et tout revient à la transitivité
\[
X\ll_{\mathrm{pol}}Y\ll_{\mathrm{pol}}Z\Longrightarrow X\ll_{\mathrm{pol}}Z.
\]
% lines 10--16
Soit $Z'\leq Z$, alors par $Y\ll_{\mathrm{pol}}Z$ \deleted{\unsure{donc}} la sous-figure induite $\unsure{Y^{4}_{Z'}|Y^{*}_{Z'}|Y^{+}_{Z'}|Y_{Z'}}$ est vide, ou de la forme $Y'$. Dans le premier cas, la figure induite par $\unsure{Y_{Z'}|Z'}$ sur $X$ est vide, OK, dans le deuxième elle est vide ou de la forme $X'$, par $\ll_{\mathrm{pol}}$, OK encore. (On utilise transitivité des images inverses)

% lines 17--22
2°) \uline{Mag 2}\quad Soit $X\ll_{\mathrm{pol}}Y$, prouvons que l'ens
\[
\widetilde{Y}^{X,\mathrm{pol}}=\{Z\in\widetilde{Y}\mid X\ll_{\mathrm{pol}}Z\}
\]
a un plus petit élément, pour ceci, il suffit de noter que l'inclusion évidente $\subset$ est une égalité
\[
\widetilde{Y}^{X,\mathrm{pol}}=\widetilde{Y}^{X},
\]
\deleted{est en fait une égalité},
% lines 23--26
i.e. que \deleted{$\unsure{X}$ \unsure{Si|$S'$}}
\[
X\ll Y'\leq Y\Longrightarrow X\ll_{\mathrm{pol}}Y',
\]
c'est évident par définition (vu que les $Y''\in\widetilde{Y}'$ sont aussi $\in\widetilde{Y}$).

% lines 27--28
\deleted{\unsure{3°)}} \uline{NB} On en conclut que
\[
X\overset{\circ}{\ll}_{\mathrm{pol}}Y\Longleftrightarrow X\ll_{\mathrm{pol}}Y\text{ et }X\overset{\circ}{\ll}Y
\]
